\section{模型的评价}

\subsection{模型的优点}

\begin{itemize}[itemsep=2pt, topsep=3pt]
	\item \textbf{物理内核统一、四问口径一致}。四个问题共用式~\eqref{eq:balance}$\sim$\eqref{eq:curtail} 的物理约束，各问只改变信息集与费用项；全部年度方案均通过逐时回代。
	\item \textbf{信息集边界严格且可自动验证}。所有预测器只使用决策时刻之前已发生的观测；五组 mutation 与计划持续性测试把“无未来信息泄漏”变成可重复执行的验收项。
	\item \textbf{风险边际有可解释依据}。$5$ 倍紧急电价在单时段下导出 $4:1$ 非对称损失和光伏输出 $0.2$ 分位启发；全年模型不直接外推该局部结论，而在多个残差分位候选中按历史窗口上的非对称费用代理在线选择。
	\item \textbf{状态与计划递推严格}。每次滚动优化以上一执行块实测 SOC 为初值；节点关闭时最近有效计划持续生效，不会回退到日初方案。
	\item \textbf{结论可归因}。问题二 A/B/C/D 拆分负载信息、风险边际与实时反馈；问题三在主节点消融外增加配对控制，分别量化负载观测刷新和光伏预报刷新；问题四区分附件 4 直用主结果与因果价格扩展。
	\item \textbf{结算与物理约束一致}。叠加结算按题面作为主口径，费用三项可直接相加；$u,v$ 等式精确表示调整方向，$z_t$ 二元变量严格排除同时充放。差额口径从头重优，仅作敏感性对照。
	\item \textbf{工程可复现}。四问均有独立求解脚本，输出保留环境信息、随机种子与结果哈希；历史参数组合可复现修订前数值。
\end{itemize}

\subsection{模型的不足}

\begin{itemize}[itemsep=2pt, topsep=3pt]
	\item 问题二与问题三采用“计划终端储电量等于当日实际初值”的日循环口径。该处理抑制日末透支，但限制跨日峰谷套利；灵敏度分析中费用影响为 \SensTerminalPct\%。
	\item 光伏预测未引入数值天气预报、辐照度、温度与云量等外生变量，仅使用同刻历史序列与残差分位数。当气象形态发生突变（例如连续阴雨转晴）时，预测与安全边际的调整会滞后一天左右。
	\item 当前安全边际对每个预测器采用全天统一的标量 kW，下修幅度未刻画日内与季节异方差；后续可按时段或季节分别估计残差分位数。
	\item 储能模型忽略自放电、容量衰减、循环寿命成本与效率随功率变化的特性，因此年度费用中没有反映电池的长期老化代价，得到的充放电深度可能比实际最优更大。
	\item 问题四按题面直接使用附件 4 全部价格；若真实市场并非提前发布完整曲线，则因果扩展更接近部署条件，但其预测精度仍可提高。
	\item 问题二采用逐日滚动的确定性等价形式求解两阶段模型，安全边际由历史残差分位数估计。当可用历史不足（如年初的前 $7$ 天）时，分位数估计的稳定性有限，本文以附件 1 的典型日曲线作为冷启动先验，但这仍是一个近似。
\end{itemize}

\section{模型的改进与推广}

\subsection{模型的改进}
\label{sec:improve}

\textbf{（1）把单日滚动扩展为多日滚动时域}。当前模型每天重新生成基准计划，跨日只传递储电量。若把优化窗口扩展到 $3\sim7$ 天并只执行首日决策，同时在窗口末端加入基于未来价格与负载的储能残值函数，就可以量化跨日套利收益并进一步降低边界效应。

\textbf{（2）把确定性等价升级为显式随机规划}。本文只在单时段获得 $0.2$ 分位启发；若显式建立负载—光伏联合情景，可进一步构造情景抽样（SAA）或条件风险价值（CVaR）模型，直接刻画跨时段风险，而无需把局部分位结论近似推广到全年。

\textbf{（3）引入更完整的价格预测}。可在电价预测中加入日内周期项、工作日／节假日效应与温度特征，或采用分位数回归直接给出价格分布，从而把电价不确定性也纳入目标函数而非仅做信息集分层。

\textbf{（4）细化储能模型}。可把循环深度相关的寿命损耗、效率曲线、自放电与维护成本写入目标函数，并用历史运行数据重新标定参数；当模型规模增大时，可采用模型预测控制与 Benders 分解等算法保持滚动求解速度。

\subsection{模型的推广}

本文的“\emph{统一物理内核 $+$ 信息集分层 $+$ 可解释风险定价}”框架不依赖本题数据。园区综合能源系统可将母线平衡扩展为电、热、气多能流平衡，并增加储热、储气状态 \cite{liu2018dayahead}；电动汽车充换电站与工业用户侧储能可分别替换负载对象、增加最大需量费；含风光互补的微网可对各电源误差分别设置非对称损失。更一般地，库存、水库与供应链等“状态逐期递推且计划—执行存在偏差结算”的问题，也可复用本文的滚动优化、单因素消融和信息集验证方法。
